\section{Threats to Validity}
\paragraph{Construct validity.} Recorded utilization is not unconstrained resource demand. ChronoGraph point samples and Alibaba completed-interval means have different semantics; external min--max values cannot be interpreted as bytes or cores. Synthetic losses affect inputs while complete recorded targets remain available. The masks simulate mechanisms rather than establish the incidence or cause of real monitoring outages. Calibration assumes the available historical prefix, not a fully missingness-degraded learning process. Availability descriptors and carry-forward summaries can retain information older than the nominal input window.

\paragraph{Internal validity.} Prefix-only normalization, partition/gap checks, future-perturbation tests for lag-MAR, fixed-mask hidden-value checks and saved-prediction audits address identifiable implementation risks. They do not certify all provider preprocessing. The external release's min--max scope and historical availability of static graph edges remain uncertain. Source validation chooses hyperparameters and stopping points; the destination prefix calibrates scale but does not fit weights. A narrow search budget, rotating neural condition exposure, full-condition analytic ridge fitting, different optimizer parameterizations and microbatch BatchNorm limit claims of equal optimization opportunity. The pretest GPU-memory amendment and recovered failures are retained, not concealed.

\paragraph{Statistical conclusion validity.} Services share a system, windows overlap and masks are paired. Seed replication estimates algorithm randomness on fixed data, not the number of independent infrastructures. The exact sign test concerns directional seed consistency, whereas the arithmetic mean difference is a separately reported effect. Source tuning reuses seed 101, addressed only by an explicitly exploratory leave-seed-out sensitivity. Conditional block-bootstrap intervals rely on a useful local block approximation; their sensitivity cannot establish universal temporal stationarity. Only four primary tests receive the prespecified multiplicity control; other patterns, horizons and cohorts are descriptive.

\paragraph{External validity.} The source and new ChronoGraph services share a period and infrastructure, the source test was historically inspected, and the destination cohort is topology-selected. Neither a local timestamped freeze nor exclusion from located earlier artifacts establishes third-party blinding. Alibaba contributes just one day, a six-hour calibration prefix and 36 target bins for 32 availability-selected services. Hundreds of millions of underlying pod records do not become hundreds of millions of independent forecast targets. Provider normalization and pod aggregation also change the task. Longer independent periods, other providers and workload regimes remain necessary before general cloud claims.

\paragraph{Efficiency and operational validity.} The laptop is shared; resident-tensor inference is not collector-to-controller latency. CPU and GPU timings depend on batch, synchronization, capture method and implementation. Parameter counts exclude state, calibration and orchestration costs. Sampled whole-system GPU memory includes other applications; process RSS is not a continuously observed peak. Epoch timing excludes some research and checkpoint costs and cannot identify a causal architecture-only training advantage. No controller, deployment, energy instrument or measured SLA benefit is available.

